% !TeX spellcheck = en_US
\documentclass[french]{yLectureNote}

\title{Électromagnétisme 3}
\subtitle{Physique}
\author{Paulhenry Saux}
\date{\today}
\yLanguage{Français}
%jesse.groenen@cemes.fr
\professor{M.Brut}
\usepackage{graphicx}%----pour mettre des images
\usepackage[utf8]{inputenc}%---encodage
\usepackage{geometry}%---pour modifier les tailles et mettre a4paper
%\usepackage{awesomebox}%---pour les boites d'exercices, de pbq et de croquis ---d\'esactiv\'e pour les TP de PC
\usepackage{tikz}%---pour deiffner + d\'ependance de chemfig
% \usepackage{tabularx}%---pour dimensionner automatiquement les tableaux avec variable X
\usepackage{awesomebox}%---Pour les boites info, danger et autres
\usepackage{menukeys}%---Pour deiffner les touches de Calculatrice
\usepackage{fancyhdr}%---pour les en-t\^ete personnalis\'ees
\usepackage{blindtext}%---pour les liens
\usepackage{hyperref}%---pour les liens (\`a mettre en dernier)
\usepackage{caption}%---pour la francisation de la l\'egende table vers Tableau
\usepackage{pifont}
\usepackage{array}%---pour les tableaux
\usepackage{lipsum}
\usepackage{yFlatTable}
\usepackage{multicol}
\usepackage{tabularx}
\usepackage{booktabs}
\newcommand{\Lim}[1]{\lim\limits_{\substack{#1}}\:}
\renewcommand{\vec}{\overrightarrow}
\newcommand{\N}[0]{\mathbb{N}}
\newcommand{\dd}{\mathrm{d}}
\newcommand{\norm}[1]{||\vec{#1}||}
\newcommand{\fo}{\psi(\vec{r},t)}
\newcommand{\foe}{\psi(\vec{r},t)\*}
\newcommand{\HH}{\hat{H}}
\newcommand{\hb}{\hbar}
\newcommand{\lap}{\nabla^2}
\newcommand{\lapcc}{\frac{\partial^2 }{\partial x^2}+\frac{\partial^2 }{\partial y^2}+\frac{\partial^2 }{\partial z^2}}
\newcommand{\E}{\vec{E}(M)}
\newcommand{\B}{\vec{B}(M)}
\newcommand{\V}{V(M)}
\newcommand{\er}{\vec{e_{\rho}}}
\newcommand{\ep}{\vec{e_{\varphi}}}
\newcommand{\pip}{\pi^+}
\newcommand{\pim}{\pi^-}
\newcommand{\mv}{\mathcal{V}}
\DeclareMathOperator{\Div}{Div}
\newcommand{\rot}{\vec{\mathrm{rot}}}
\newcommand{\grad}{\vec{\mathrm{grad}}}
%\setcounter{chapter}{3}
\begin{document}
%\chapter{Puissance et énergie du champ Em}
\setcounter{chapter}{4}
\chapter{Ondes électromahnétiques dans un milieu matériel}
\section{Équations de Maxwell}
\subsection{Équations de Maxwell dans un milieu matériel}
On  a dans un milieu matériel
\begin{theorem}[Équations de Maxwell dans un milieu matériel]
    \begin{itemize}
    \item $\rot \vec{H} - \frac{\partial \vec{D}}{\partial t}=\vec{j_{libre}}$
    \item $\rot \vec{E}+\frac{\partial \vec{B}}{\partial t} = \vec{0}$
    \item $\Div \vec{D} = \rho_{libre}+\rho_{autres}$
    \item $\Div \vec{B} = 0$
\end{itemize}
\end{theorem}

\begin{theorem}[Relations de passage]
    \begin{flalign}
& (\vec{D}_2 - \vec{D}_1) \cdot \vec{n}_{1\to2} = \sigma && \text{(Composante normale de } \vec{D}\text{)} & \\
& (\vec{E}_2 - \vec{E}_1) \times \vec{n}_{1\to2} = \vec{0} && \text{(Composante tangentielle de } \vec{E}\text{)} & \\
& (\vec{B}_2 - \vec{B}_1) \cdot \vec{n}_{1\to2} = 0 && \text{(Composante normale de } \vec{B}\text{)} & \\
& (\vec{H}_2 - \vec{H}_1) \times \vec{n}_{1\to2} = \vec{j}_S && \text{(Composante tangentielle de } \vec{H}\text{)} &
\end{flalign}
\end{theorem}
\section{Grandeurs énergétiques}
Dans le vide :
\begin{itemize}
    \item Densité volumique dénergie em : $w = \frac{\epsilon_0 E^2}{2}+ \frac{B^2}{2\mu_0}$
    \item Vecteur de Poyting : $\Pi = \vec{E}\wedge \frac{\vec{B}}{\mu_0}$ défini à partir du champ des réels 
    \item Pour une onde progressive dans la direction u : $\vec{Pi} = wc\vec{u}$ : Pi est un courant volumique d'énergie em.
\end{itemize}

Dans le milieu matériel linéaire :
\begin{itemize}
    \item Densité volumique dénergie em : $w = \frac{\vec{D}\cdot \vec{E}}{2}+\frac{\vec{H}\cdot \vec{B}}{2}$
    \item Vecteur de Poyting : $\Pi = \vec{E}\wedge \vec{H}$ défini à partir du champ des réels 
\end{itemize}
On introduit la puissance moyenne rayonnée à travers une surface S :
$$P = \iint_S <\vec{\Pi}>_{temps}\cdot \vec{\dd S}$$
L'intensité vaut alors $ I = \frac{P}{S}$
\section{Ondes électromagnétiqus dans un milieu Diélectrique}
\subsection{Ondes monochormatiques dans un milieu diélectrique LHI}
On consdidère des milueux localment neutres, isolants et non magnétique :
\begin{itemize}
    \item $\rho_{libre} =0$
    \item $\vec{j_{libre}} = 0$
    \item $\vec{M} = \vec{0}$
    \item $\vec{B} = \mu_0\vec{H}$
\end{itemize}
On introduit la relation constitutive du milieu : $\vec{\underline{D}} = \epsilon_0\underline{\epsilon_r}(\omega)\vec{\underline{E}}$ et $\epsilon_r = 1+ \chi_e(\omega)$

Les équations de Maxwell sont identiques à l'exception que $\epsilon_0\to \epsilon_0\epsilon_r$ :Les ondes em qui peuvent se propager dans le milieu sont semblables à celles de propageant dans le vide.

$$\lap \vec{\underline{E}}+\frac{\omega^2}{c^2}\epsilon_r(\omega)\vec{\underline{E}}=\vec{0}$$
avec la vitesse de propagation : $\frac{c}{\sqrt{\epsilon_r}}$

\subsection{Ondes planes progressives monochromatiques (OPPM)}

Une onde plane se caractérise par des amplitudes identiques en tout point d'un plan perpendiculaire à la direction de propagation. On se place dans un milieu linéaire, homogène et isotrope (LHI).

\begin{definition}{Structure de l'OPPM dans un milieu LHI}
En recherchant des solutions sous forme de vagues planes complexes de la forme $\vec{\underline{E}} = \vec{\underline{E}}_0 \exp\left(i(\omega t - \vec{k}\cdot\vec{r})\right)$, les équations de Maxwell se réduisent aux relations algébriques suivantes :
\begin{itemize}
    \item Transversalité : $\vec{k}\cdot \vec{\underline{E}} = 0 \quad \text{et} \quad \vec{k}\cdot \vec{\underline{B}} = 0$
    \item Relation de structure : $\vec{k}\wedge \vec{\underline{E}} = \omega \vec{\underline{B}}$
    \item Équation d'Ampère-Maxwell : $\vec{k}\wedge \vec{\underline{B}} = -\frac{\omega}{c^2}\epsilon_r(\omega)\vec{\underline{E}}$
\end{itemize}
\end{definition}

\tipsInfo{Trièdre direct}{La relation de structure implique que le trièdre $(\vec{k}, \vec{\underline{E}}, \vec{\underline{B}})$ est direct. Les champs électrique et magnétique vibrent en phase dans un milieu non absorbant.}

L'équation de propagation (équation de Helmholtz) pour le champ électrique se simplifie sous la forme :
$$k^2 \vec{\underline{E}} = \frac{\omega^2}{c^2}\epsilon_r(\omega)\vec{\underline{E}}$$

La nature de la permittivité diélectrique relative $\epsilon_r(\omega)$ détermine le comportement de l'onde dans le milieu :

\begin{proposition}{Propriétés du milieu selon $\epsilon_r(\omega)$}
\begin{itemize}
    \item \textbf{Milieu transparent} : $\epsilon_r(\omega)$ est réel et strictement positif. Le vecteur d'onde $k = \frac{\omega}{c}\sqrt{\epsilon_r(\omega)}$ est réel : l'onde se propage sans atténuation de son amplitude.
    \item \textbf{Milieu dispersif} : $\epsilon_r$ dépend explicitement de la pulsation $\omega$. Des ondes de fréquences différentes se propagent à des vitesses différentes.
    \item \textbf{Milieu non dispersif} : $\epsilon_r = \text{cste}$. C'est le cas unique du vide où $\epsilon_r = 1$.
\end{itemize}
\end{proposition}

\warningInfo{Cas des milieux conducteurs ou absorbants}{Si $\epsilon_r(\omega)$ possède une composante imaginaire non nulle, le vecteur d'onde $\vec{k}$ devient complexe ($\vec{\underline{k}} = \vec{k}' - i\vec{k}''$), ce qui traduit physiquement une atténuation (effet Joule ou absorption) de l'onde lors de sa propagation.}

On introduit l'indice de réfraction du milieu $n(\omega) = \sqrt{\epsilon_r(\omega)}$, ainsi que la vitesse de phase de l'onde $v_{\varphi}(\omega)$ :

$$v_{\varphi}(\omega) = \frac{\omega}{k} = \frac{c}{n(\omega)}$$

\checkInfo{Longueur d'onde}{Dans un milieu matériel d'indice $n(\omega)$, la longueur d'onde est réduite par rapport à sa valeur dans le vide $\lambda_0$ :
$$\lambda = \frac{2\pi}{k} = \frac{\lambda_0}{n(\omega)}$$
}

\subsection{Aspects énergétiques}

Le transport d'énergie par l'onde électromagnétique est décrit par le vecteur de Poynting $\vec{\Pi}$. 

\criticalInfo{Calcul de la puissance instantanée}{Pour les aspects énergétiques (produits de champs), il est impératif d'utiliser les \textbf{champs réels} et non les expressions complexes ! 
$$\vec{\Pi} = \frac{\vec{E}_{\text{réel}} \wedge \vec{B}_{\text{réel}}}{\mu_0}$$
}

Pour une OPPM se propageant selon le vecteur unitaire $\vec{u}$ dans un milieu transparent d'indice $n(\omega)$, le vecteur de Poynting instantané s'écrit :
$$\vec{\Pi} = n(\omega)c\epsilon_0 E_{\text{réel}}^2\vec{u}$$

\tipsInfo{Moyenne temporelle}{En pratique, les détecteurs (œil, photodiodes) sont sensibles à la puissance moyenne. Pour une OPPM, la valeur moyenne temporelle du vecteur de Poynting (l'éclairement) vaut :
$$\langle \vec{\Pi} \rangle = \frac{1}{2}n(\omega)c\epsilon_0 E_0^2\vec{u}$$
}

\subsection{Milieux absorbants}
On a $\epsilon_r$ complexe donc $\vec{k}$ complexe.

\begin{itemize}
    \item Si $\underline{\epsilon_r} = \epsilon_r'$ : 
\begin{itemize}
    \item Si $\epsilon_r' >0$ : propagation sans atténuation : milieu transparent : $e^{i(kz-\omega t)}$
    \item Si $\epsilon_r'<0$ : pas de propagation, que de l'atténuation : $e^{-kz}e^{-i\omega t}$ 
\end{itemize}
\item Dans le cas contraire,  absorption, propagation et atténuation
\end{itemize}

\subsection{Réflexion et transmission d'une onde à l'interface entre 2 diélectriques}

\includegraphics[scale=0.5]{resume.pdf}
\end{document}
